
\medskip

\begin{enumerate}
	\item \begin{tabularx}{\linewidth}[t]{@{}X@{\quad}l}
En suivant le programme, une fois que le stylo est en position d'écriture, le chat est orienté vers la droite, et il avance d'abord de 100 pas, d'où un premier segment horizontal de 100 pas (donc 5 carreaux sur la figure).
		
Ensuite, il tourne vers la droite, de \textbf{angle}°, soit ici 90°, donc il est maintenant orienté vers le bas. 
		
C'est là qu'il avance de \textbf{longueur} pas, donc, ici de 80 pas (soit un segment vertical de 4 carreaux sur la figure).
		
Puis il tourne de $180 - \text{\textbf{angle}} = 180 - 90 = 90$°. Il est donc maintenant orienté vers la gauche.
		
Puis, le programme continue pour terminer le parallèlogramme, qui, ici, sera un rectangle. On obtient la figure ci-contre:&
		
\begin{tikzpicture}[x=0.5cm,y=0.5cm,baseline={(0,5)}]
\draw [gray!70!, xstep=1,ystep=1] (0.2,0.2) grid (10.8,10.8);
\draw [line width=1pt] (0,0) rectangle (11,11);
\draw [line width=1.5pt] (2,9) rectangle (7,5);
\node at (2,9) {\includegraphics[width=0.7cm]{chat}};
%\node at (5,0) [below=2mm]{Le côté d'un carreau représente 20 pas};
\end{tikzpicture}
	\end{tabularx}
	\item Puisque le parallélogramme est aussi un rectangle dans ce deuxième exemple, on va également choisir \textbf{angle} = 90.
	
	Ce parallélogramme étant même un carré, il faut que les quatre côtés soient de même longueur, de 100 pas sur la figure, donc on choisit \textbf{longueur} = 100.
	
	\item Sur les trois figures proposées la longueur des côtés qui ne sont pas horizontaux est bien la moitié de la longueur des côtés horizontaux, donc seul l'angle permettra de trancher. 
	
Avec le rappel fait au début de l'exercice, si l'angle saisit est de $75$\degres, l'angle constaté entre les deux segments tracés sera l'angle supplémentaire, dont la mesure sera donc : $ 180 - 75 = 105 $\degres.
\end{enumerate}


